%% Section 3 — Results
%% CMPB submission: TRUST-CT
%% All values drawn from locked artifacts; see numerical_provenance.csv.
%% DO NOT edit numerical values without updating numerical_provenance.csv.

\section{Results}
\label{sec:results}

\subsection{Cohort and run completeness}
\label{sec:cohort_completeness}

Cimas comprised 12{,}649 participants across 23 providers;
BETTER-BP comprised 402 randomized participants across two providers;
eICU comprised 2{,}520 stays across eight hospital groups.
For the eICU security benchmark, 60 features were selected by
training-only mutual information (SHA-256: \texttt{4f4a6603\ldots});
the selected list and full hyperparameter configuration are
reported in Supplementary~S1.
Of 1{,}200 intended Phase~6 experimental outcomes, all 1{,}200 were
reconciled (1{,}065 valid-primary runs, 35 Krum-inapplicable
diagnostic fallbacks, and 100 no-attack f-invariant skips);
zero unexplained failures occurred (Supplementary~S1).

\subsection{Federated predictive performance}
\label{sec:fl_results}

% W03b fix: separated rounded-marginal difference from bootstrap NI estimate.
% $\Delta=-0.004$ was from rounded AUROC marginals (0.795-0.799);
% the provider-paired bootstrap estimate is $\hat{d}=-0.002$ (obs=-0.002135).
% LCB=-0.0045 refers to the bootstrap one-sided lower bound, not the
% marginal difference. These must not be conflated.
All five Cimas FL methods were non-inferior to centralized training
on the prespecified AUROC margin of $-0.02$
(Table~\ref{tab:prediction}; Figure~\ref{fig:ni}).
FedAvg achieved estimated AUROC 0.795 [95\%~CI 0.773--0.811] versus
centralized 0.799 [0.778--0.814].
The provider-paired bootstrap NI estimate was
$\hat{d}=-0.002$ (95\%~CI $-0.005$, $+0.001$;
one-sided NI lower confidence bound $-0.0045 > -0.02$;
\emph{non-inferior}).
The worst NI lower confidence bound across all methods was $-0.0095$
(FedAdam; non-inferior; upper CI $+0.000$).
FedAvg AUROC exceeded local-only training by $+0.007$
[95\%~CI $-0.003$, $+0.018$], but the interval included zero.

% W02c fix: N=3716 long-tail are held-out providers from the SAME Cimas
% data source — NOT an external validation set. Relabelled accordingly.
On the long-tail within-source transportability evaluation set ($N=3{,}716$;
non-adherence prevalence 0.677), the centralized model achieved
risk-positive PR-AUC 0.900 [0.882, 0.917] (PR-lift 1.329);
FedAvg achieved PR-AUC 0.899 [0.880, 0.916] (PR-lift 1.327).
These results are from the Cimas LOCO evaluation and are distinct from
the eICU security benchmark reported in Section~\ref{sec:adversarial}.

\subsection{Policy evaluation and closed-loop replay}
\label{sec:policy_results}

The DR-estimated value of the learned uplift-targeting policy was
$\hat{V}(\pi_\text{up})=0.892$ versus the random-budget comparator
$\hat{V}(\pi_\text{rand})=0.860$
($\Delta=+0.032$; 95\%~CI $[-0.011, +0.078]$;
Table~\ref{tab:prediction}).
The confidence interval crossed zero; no statistically supported
policy-value improvement was established.
The incentive improved average estimated attendance relative to
no-incentive ($\hat{V}(\pi_0)=0.815$), but available
baseline predictors could not identify differential responders.

The Cimas closed-loop replay reproduced
$\hat{V}(C1, B=\tfrac{2}{3})=0.8603$ from calibration (target 0.860).
BETTER-BP simulated participation was:
C0 81.4\% [81.3, 81.5\%], C1 85.9\% [85.7, 86.0\%],
C2 87.2\% [87.1, 87.3\%], C3 89.0\% [88.9, 89.1\%].
The Cimas cold-start nAULC differences (C2$-$C1 and C3$-$C1)
were approximately $7.8\times10^{-5}$; confidence intervals
may exclude zero but absolute differences remained below the
prespecified operational materiality threshold of 0.001
(operationally negligible; Figure~\ref{fig:policy}).

\subsection{Adversarial robustness}
\label{sec:adversarial}

The eICU clean FedAvg AUROC was 0.8904 [0.8892, 0.8913]
across five seeds and 15 rounds.
Table~\ref{tab:security} and Figure~\ref{fig:attacks} summarise
attack and defense results.

At $f=3$ (37.5\% malicious clients), label flipping produced
the largest aggregate AUROC degradation
($\Delta=-0.120$ [$-0.178$, $-0.061$]; $p=0.005$).
Sign flipping caused $\Delta=-0.059$ [$-0.103$, $-0.016$]; $p=0.020$.
Model replacement and ALIE caused smaller reductions
($\Delta=-0.020$ and $\Delta=-0.004$, respectively;
ALIE: $p=0.24$).
% W10a fix: BSR renamed; denominator clarified; clean-model baseline absent.
% W09c fix 2026-08-28: original backdoor used fixed seed=42 (same poisoning
% indices across all outer seeds). Corrected to seed from (outer_seed,round,client).
% BSR changed from 0.922 to 0.920 (difference 0.002, within noise; conclusion unchanged).
The targeted tabular backdoor achieved a triggered target-class rate
(proportion of all triggered test rows predicted as target class 0) of 0.920.
No clean-model baseline was collected; this metric cannot be attributed to the
backdoor without a clean-model comparison and is therefore reported descriptively.
Aggregate AUROC was not materially reduced ($+0.001$; $-0.001$, $+0.002$),
demonstrating that AUROC alone did not detect the targeted attack.

Coordinate-wise median aggregation recovered
$+0.111$ [0.058, 0.165] AUROC relative to FedAvg under
label flipping ($p=0.005$); the worst residual degradation
under coordinate median across all attacks was approximately $-0.008$.
Krum was not applicable at $f=3$ for eICU ($K=8$; $8\not>8$);
the 35 fallback rows are excluded from Krum estimates.

\subsection{Privacy--utility evaluation}
\label{sec:privacy}

% W09b fix 2026-08-28: original 6C-i noise seeded seed*1000+t gave all clients
% same noise per round (same vector). Corrected to per-client seed
% outer_seed*1_000_000 + round_t*1_000 + client_idx. This inflated utility
% degradation at alpha>=0.01 in original; corrected values below.
At the selected operating point $\alpha=0.01$, the utility cost was
$\Delta\text{AUROC}=-0.001$ [$-0.002$, $+0.001$] ($p=0.26$),
remaining well within the prespecified materiality range
(Table~\ref{tab:security}; Figure~\ref{fig:frontier}).
At $\alpha=0.05$, the corrected per-client noise degradation was
$\Delta\text{AUROC}=-0.011$ [$-0.019$, $-0.000$] ($p=0.09$),
also within the materiality threshold of $-0.02$
(the original value $-0.074$ was inflated by the W09b shared-noise bug).
The analytic single-step gradient canary confirmed exact feature
recovery (cosine similarity 1.000000), establishing the theoretical
upper bound for direct gradient access.
Under the coordinator-observed multi-step FedAvg delta, the same
approximation formula yielded cosine $-0.538$, confirming the
formula is invalid for multi-epoch updates.
% W07 resolved 2026-08-28: C0-correct batch=1 reconstruction cosines confirmed
% from trust-ct/processed/phase6/phase6c_reconstruction.csv (125 batch=1 rows;
% 5 hospitals x 5 rounds x 5 alpha values; C0-correct MI features, seed=7 split).
% Prior stale values (0.534/0.464 from audit4_leakage_table.csv) were from
% pre-C0 MI feature set pooled across batch sizes 1/4/16 — not comparable.
% Feature-sign recovery was not stored per-sample in phase6c_reconstruction.csv.
Under a single-step fresh-gradient reconstruction at batch size~1,
cosine similarity was 1.000 ($\alpha=0$; exact canary recovery),
$0.996\pm0.003$ ($\alpha=0.01$), and $0.879\pm0.070$ ($\alpha=0.05$).
Reconstruction degrades sharply with batch size
(cosine $0.307$ at batch~4 and $0.009$ at batch~32, $\alpha=0$).
Feature-sign recovery at batch size~1 was not separately stored in the
C0-correct reconstruction run;
\emph{[AUTHOR\_PENDING: sign recovery requires re-extraction from stored
reconstruction vectors --- see trust-ct/remediation/w07\_batch1\_corrected\_values.json]}.
% W09b stale alpha=0.01 paragraph removed 2026-08-29. Correct value
% ($-0.001$ [$-0.002$, $+0.001$], $p=0.26$) stated at Section 3.5 line 1 above
% and in tab:security. Do not reintroduce the old $-0.005$ figure.

MI attack AUROC was 0.497 [0.496, 0.498] (eICU; Cohen's $h=-0.006$)
and 0.507 [0.506, 0.507] (Cimas; Cohen's $h=+0.014$).
The eICU attacker performed at or below chance.
The Cimas attacker showed a $+0.7$ percentage-point AUROC deviation;
the effect size was operationally negligible ($|h|<0.02$).
The near-chance eICU baseline (AUROC 0.497 at $\alpha=0$)
indicates that model confidence was insufficient to meaningfully
separate members from non-members prior to any perturbation;
the improvement at $\alpha=0.01$ cannot be attributed to update noise.

\subsection{Audit-log evaluation}
\label{sec:audit}

% W11b fix: tests verified hash-link integrity, NOT signature non-repudiation.
% Removed "signed" from first sentence to avoid conflating HMAC auth with
% public-key signature.
The hash-linked audit log detected all 350 prespecified injected events
(350/350; zero false alerts) across seven evaluated fault types
(modification, replay, reordering, stale attestation, update substitution,
dropout, and silent dropout), ten injection positions, and five seeds
(Table~\ref{tab:ablation}).
These tests evaluated hash-chain integrity under the specified mutation
scenarios; they did not evaluate public-key signature non-repudiation,
which the HMAC-based mechanism does not provide.
The log was not evaluated for faults outside the seven prespecified
categories or for events submitted by authorized sources with
incorrect content.

\subsection{Component ablation}
\label{sec:ablation}

Table~\ref{tab:ablation} summarises five component ablations (A1--A5).
Federated learning was non-inferior to centralized training (A1).
Coordinate-wise median aggregation materially mitigated label-flip
poisoning (A2; $\Delta=+0.111$, $p=0.005$).
The $\alpha=0.01$ perturbation remained within the prespecified
utility materiality range while reducing the optimistic
reconstruction signal (A3; $\Delta\text{AUROC}=-0.001$ [$-0.002$, $+0.001$], $p=0.26$).
The audit log detected all evaluated governance faults with
zero false alerts (A4).
DR-policy uplift was not established (A5; CI crosses zero).

\subsection{Provider-level performance disparity: warm/cold de-confound}
\label{sec:e6}

% E6 experiment: 4 conditions x 5 seeds x 60 rounds, Cimas K=23.
% Data: processed/e6_fairness/e6_provider_round_results.csv (20,700 rows)
% Analysis: run_e6_analysis.py, run_e6_disparity.py
% warm_fed exactly reproduces cold_fed rounds 31-60 (3,450/3,450 identical;
% max |delta|=0); it is a reproducibility gate, not an independent arm.
% Provenance: e6_provider_fixed_effects.csv, e6_disparity_stats.csv,
%             e6_paired_analysis.csv, e6_analysis_report.txt

To isolate the mechanism of federation benefit, a four-condition experiment
was conducted on the Cimas primary cohort
($K=23$ providers, five seeds, 60~rounds; Supplementary~S5).
Conditions: \emph{cold-FedAvg} (FedAvg from zero); \emph{cold-local}
(independent local training from zero); \emph{warm-local} (local-only from
the round-30 FedAvg checkpoint); and a reproducibility gate (\emph{warm-FedAvg},
continuing FedAvg from the same checkpoint, which numerically reproduced
cold-FedAvg rounds~31--60 exactly: 3,450/3,450 pairs; max $|\delta|=0$).
All conditions shared identical test sets, preprocessing, features,
local-epoch count ($N_\text{local}=5$), and learning rate ($\eta=0.02$).
Performance was averaged over rounds~51--60 per (provider, seed).
Inference: seed-level paired $t$-tests ($n=5$, df=4) with $5{,}000$-resample
seed-cluster bootstrap 95\%~CIs.

\paragraph{Condition-level provider performance disparity.}
The following metrics describe AUROC dispersion across the $K=23$~Cimas providers;
they measure provider-level performance equity, not demographic or individual fairness.
Under cold FedAvg the provider AUROC interquartile range (IQR) was
0.092 [0.073, 0.109] versus 0.094 [0.073, 0.116] under cold local;
the distributions were comparable.
Warm-local achieved a slightly narrower IQR (0.091 [0.075, 0.105]).
Worst-provider absolute AUROC was 0.518 [0.444, 0.578] (cold FedAvg),
0.509 [0.466, 0.566] (cold local), and 0.525 [0.487, 0.565] (warm-local).

\paragraph{Paired contrasts (macro-average, late rounds 51--60).}
Cold FedAvg versus cold local:
macro-average $\Delta=+0.003$ [$-0.003$, $+0.008$]; $p=0.39$;
provider-weighted $\Delta=+0.002$; $p=0.14$.
No significant bulk AUROC benefit was established.
The worst-provider contrast delta was $-0.114$ [$-0.141$, $-0.089$]; $p=0.001$:
the single most data-heterogeneous provider was materially harmed by
global aggregation.
Zero of 23~providers showed systematic directional effects
(signal-to-noise ratio $|\bar{\delta}|/\mathrm{SD}_\delta > 2$);
the effect was highly heterogeneous (SD of provider mean deltas: 0.032~AUROC).

Warm-local versus cold local:
provider-weighted $\Delta=+0.002$ [$+0.001$, $+0.003$]; $p=0.012$.
The round-30 federated checkpoint provided the dominant component
of the bulk FedAvg benefit with no further aggregation required.

Warm-local versus continued FedAvg:
macro $\Delta=+0.002$ [$+0.000$, $+0.006$]; $p=0.21$;
provider-weighted $\Delta=+0.000$; $p=0.86$.
The two conditions were statistically indistinguishable in bulk AUROC.
Warm-local required a single weight broadcast
($529$~parameter floats across $23$~providers, $22$~features plus intercept)
versus 30~aggregation rounds for continued FedAvg ($31{,}740$~floats),
a 98.3\% reduction in parameter transfers with no significant average AUROC cost.

These results indicate that shared weight initialisation, not sustained
gradient aggregation, accounts for the primary performance parity in this
setting. Provider-level contrast heterogeneity substantially exceeds the
macro-average effect, indicating that aggregate summaries mask meaningful
provider-specific variation.

%% ---- Tables ----

\begin{table}[ht]
\centering
\caption{Federated prediction (Cimas LOCO, $N=12{,}649$, $K=23$)
  and policy evaluation (BETTER-BP, $N=402$, $B=\tfrac{2}{3}$).}
\label{tab:prediction}
\begin{tabular}{@{}llll p{2.5cm} p{2.5cm}@{}}
\hline
Analysis & Comparator & Metric & Estimate [95\% CI] & Criterion & Interpretation \\
\hline
% W03b corrected 2026-08-29: -0.004 was rounded AUROC marginal (0.795-0.799);
% provider-paired bootstrap estimate is -0.002 (obs=-0.002135 from ni_client_level.csv).
FedAvg vs Central   & Centralized      & $\Delta$AUROC & $-0.002$ [$-0.005$, $+0.001$] & NI LCB $>-0.02$ & Non-inferior \\
FedAdam vs Central  & Centralized      & $\Delta$AUROC & $-0.005$ [$-0.011$, $0.000$]  & NI LCB $>-0.02$ & Non-inferior \\
Central (long-tail) & Prevalence 0.677 & PR-AUC        & $0.900$ [0.882, 0.917]        & PR-lift $>1$    & Lift $=1.329$ \\
FedAvg (long-tail)  & Prevalence 0.677 & PR-AUC        & $0.899$ [0.880, 0.916]        & PR-lift $>1$    & Lift $=1.327$ \\
\hline
$\pi_\text{up}$ vs $\pi_\text{rand}$ & Random budget & $\Delta\hat{V}$ & $+0.032$ [$-0.011$, $+0.078$] & CI excludes 0 & Not established \\
\hline
\end{tabular}
\end{table}

\begin{table}[ht]
\centering
\caption{Adversarial robustness and privacy--utility evaluation
  (eICU, $K=8$, $f=3$, 5~seeds, 15~rounds).}
\label{tab:security}
\begin{tabular}{@{}lllll@{}}
\hline
Condition & Aggregator / $\alpha$ & Utility change ($\Delta$AUROC) & Leakage metric & Interpretation \\
\hline
Clean baseline     & FedAvg        & ---                              & AUROC 0.890 [0.889, 0.891]        & Reference \\
Label flip ($f=3$) & FedAvg        & $-0.120$ [$-0.178$, $-0.061$]$^\ast$ & BSR n/a; PR-AUC $\downarrow$ & Poisoning effective \\
Label flip ($f=3$) & Coord. median & $+0.111$ [$+0.058$, $+0.165$]$^\ast$& Recovery vs FedAvg          & Defense effective \\
Sign flip ($f=3$)  & FedAvg        & $-0.059$ [$-0.103$, $-0.016$]$^\dagger$ & ---                    & Poisoning effective \\
Backdoor ($f=3$)   & FedAvg        & $+0.001$ [$-0.001$, $+0.002$]  & BSR 0.920                         & AUROC insensitive \\
ALIE ($f=3$)       & FedAvg        & $-0.004$ [$-0.013$, $+0.004$]  & ---                               & Not significant \\
\hline
% W07+W09b resolved 2026-08-28: cosines from C0-correct recon; noise from per-client W09b.
% Sign recovery not stored per-sample; remains AUTHOR_PENDING.
$\alpha=0.01$ & Clipping$+$noise & $-0.001$ [$-0.002$, $+0.001$]$^{\ddagger}$ & Recon cos.\ $0.996\pm0.003$ (batch=1) & Within materiality \\
$\alpha=0.05$ & Clipping$+$noise & $-0.011$ [$-0.019$, $-0.000$] & Recon cos.\ $0.879\pm0.070$; sign rec.\ [AUTHOR\_PENDING] & Within materiality (W09b) \\
eICU MI attack  & $\alpha=0$    & ---                              & AUROC 0.497 [0.496, 0.498]; $h=-0.006$ & At/below chance \\
Cimas MI attack & $\alpha=0$    & ---                              & AUROC 0.507 [0.506, 0.507]; $h=+0.014$ & Negligible \\
\hline
\end{tabular}
\begin{tablenotes}
\small $^\ast$ Paired $t$-test, $p=0.005$; $^\dagger$ $p=0.020$;
% W05a fix: 500 resamples confirmed from N_BOOT in run_phase5_v3.py.
$^\ddagger$ $p=0.39$; 95\% CIs from bootstrap (500 resamples).
\end{tablenotes}
\end{table}

\begin{table}[ht]
\centering
\caption{Component ablation summary (A1--A5).}
\label{tab:ablation}
\begin{tabular}{@{}lllll@{}}
\hline
ID & Component & Locked estimate [95\% CI] & Prespecified criterion & Conclusion \\
\hline
A1 & FL noninferiority (FedAvg)      & NI LCB $=-0.0045$                 & $>-0.02$ AUROC            & Supported \\
A2 & Coord.\ median under label flip  & Recovery $+0.111$ [0.058, 0.165]  & $p<0.05$; $\Delta>0$      & Supported \\
A3 & $\alpha=0.01$ utility cost       & $\Delta=-0.001$ [$-0.002$, $+0.001$] & $|\Delta|<0.02$; $p=0.26$ & Within range \\
A4 & Audit-log fault detection        & 350/350 (0 false alerts)          & All 350 detected          & Supported \\
A5 & DR-policy uplift                 & $+0.032$ [$-0.011$, $+0.078$]    & CI excludes 0             & Not established \\
\hline
\end{tabular}
\end{table}
